% (c) Nikita Lisitsa, lisyarus@gmail.com, 2026

\documentclass[10pt]{beamer}

\usepackage[T2A]{fontenc}
\usepackage[russian]{babel}
\usepackage{minted}

\usepackage[most]{tcolorbox}
\tcbuselibrary{minted}

\usepackage{graphicx}
\graphicspath{ {./images/} }

\usetheme{metropolis}

\usepackage{adjustbox}

\usepackage{color}
\usepackage{soul}

\usepackage{hyperref}

\definecolor{codebg}{RGB}{29,35,49}

\setminted{fontsize=\footnotesize}

\setlength\partopsep{-\topsep}

\definecolor{blue}{rgb}{0,0,1}
\definecolor{red}{rgb}{1,0,0}

\usemintedstyle{tango}

\makeatletter
\newcommand{\slideimage}[1]{
  \begin{figure}
    \begin{adjustbox}{width=\textwidth, totalheight=\textheight-2\baselineskip-2\baselineskip,keepaspectratio}
      \includegraphics{#1}
    \end{adjustbox}
  \end{figure}
}
\makeatother

\title{Компьютерная графика}
\subtitle{Практика 7: освещение}
\date{2026}

\setbeamertemplate{footline}[frame number]

\begin{document}

\frame{\titlepage}

\begin{frame}[fragile]
\frametitle{Практика 7}
\begin{itemize}
\item Камеру можно двигать и вращать стрелками
\end{itemize}
\end{frame}

\begin{frame}[fragile]
\frametitle{Практика 7}
\slideimage{0.png}
\end{frame}

\begin{frame}[fragile]
\frametitle{Задание 1}
Добавляем направленный источник света
\begin{itemize}
\item Добавляем в \mintinline{cpp}|LightsUniforms| направление на источник света (нормированный вектор!) и его яркость
\item Направление должно как-то меняться со временем
\item \textbf{NB:} \mintinline{wgsl}|vec3f| в uniform-буфере должен быть выровнен на 16 байт! Проще всего этого добиться, добавив руками padding-байты
\item Добавляем в \mintinline{cpp}|Immediates| цвет (альбедо) нашего объекта (вектор с компонентами от 0 до 1)
\item Во фрагментном шейдере (в \mintinline{text}|shader.wgsl|) меняем вычисление цвета на формулу диффузного освещения (без нормализации)
\end{itemize}
\usemintedstyle{lightbulb}
\begin{tcblisting}{listing engine=minted, minted language=wgsl, minted options={fontsize=\scriptsize}, listing only, colback=codebg, colframe=codebg, rounded corners}
albedo * lightIntensity
  * (ambient + max(0.0, dot(normal, lightDirection)))
\end{tcblisting}
\end{frame}

\begin{frame}[fragile]
\frametitle{Задание 1}
\slideimage{1.png}
\end{frame}

\begin{frame}[fragile]
\frametitle{Задание 2}
Визуализируем точечный источник света
\begin{itemize}
\item Заводим переменную под положение точечного источника (можно сразу в \mintinline{cpp}|LightsUniforms|)
\item Он тоже должен меняться со временем
\item Визуализируем его с помощью \mintinline{cpp}|lightRenderPipeline|: ей нужно выставить \mintinline{cpp}|LightImmediates| с позицией источника света, и нарисовать с помощью \mintinline{cpp}|render_pass.draw(vertex_count = 6)|
\end{itemize}
\end{frame}

\begin{frame}[fragile]
\frametitle{Задание 2}
\slideimage{2.png}
\end{frame}

\begin{frame}[fragile]
\frametitle{Задание 3}
Добавляем освещение от точечного источника
\begin{itemize}
\item Добавляем яркость и параметры затухания в \mintinline{cpp}|LightsUniforms|
\item Коэффициенты затухания -- тоже 3D вектор, например \((1, 0, 0.1)\)
\item Добавляем освещение от источника в шейдер
\item Направление на источник нужно вычислить на основе его положения и положения текущего пикселя в пространстве \mintinline{cpp}|in.world_position|
\item Затухание нужно вычислить на основе расстояния до источника (см. лекцию)
\item \textbf{NB:} здесь удобно выделить функцию \mintinline{cpp}|computeLighting|, обрабатывающую один источник, зная его направление и цвет
\end{itemize}
\end{frame}

\begin{frame}[fragile]
\frametitle{Задание 3}
\slideimage{3.png}
\end{frame}

\begin{frame}[fragile]
\frametitle{Задание 4}
Добавляем specular освещение
\begin{itemize}
\item Добавляем в \mintinline{cpp}|Immediates| параметры \mintinline{cpp}|roughness| (например, 0) и \mintinline{cpp}|specularFactor| (например, 1)
\item \mintinline{cpp}|specularFactor| -- просто множитель для specular слагаемого
\item На основе \mintinline{cpp}|roughness| вычисляем \mintinline{cpp}|power| для модели Блинна-Фонга (с \mintinline{cpp}|halfway| вектором)
\item Вычисляем направление из пикселя в сторону камеры, и на его основе -- \mintinline{cpp}|halfway| вектор
\item Используя их, вычисляем specular компоненту \textbf{(на этот раз с нормализацией!)}
\end{itemize}
\end{frame}

\begin{frame}[fragile]
\frametitle{Задание 4}
\slideimage{4.png}
\end{frame}

\begin{frame}[fragile]
\frametitle{Задание 5*}
Рисуем 9 копий модели
\begin{itemize}
\item У каждой копии должны быть свои значения \mintinline{cpp}|roughness| и \mintinline{cpp}|specularFactor|
\item У каждой копии должна быть своя матрица \mintinline{cpp}|model|
\item Должно быть достаточно 9 раз выставить \mintinline{cpp}|Immediates| и сделать \mintinline{cpp}|render_pass.draw_indexed()|
\end{itemize}
\end{frame}

\begin{frame}[fragile]
\frametitle{Задание 5*}
\slideimage{5.png}
\end{frame}

\end{document}
